\documentclass[10pt,a4paper]{amsart}
\usepackage[margin=19mm]{geometry}
\usepackage{amsmath,amssymb,mathtools}
\usepackage[scr=rsfs,cal=euler]{mathalfa}
\usepackage{xfrac}
\usepackage[unicode,hidelinks,bookmarks=false]{hyperref}
\newcommand{\ie}{i.e.\ }
\newcommand{\cf}{cf.\ }
\newcommand{\resp}{respectively\ }
\newcommand{\Subsection}{\subsection}
\providecommand{\nobreakdash}{\nobreak\hbox{-}}
\setlength{\emergencystretch}{2em}
\multlinegap0pt
\usepackage{bm}
\usepackage{bbm}
\usepackage{dsfont}
\usepackage{xspace}
\usepackage{cancel}
\usepackage{mathtools}
\usepackage{amsthm}
\makeatletter
\@ifundefined{definition}{\newtheorem{definition}{Definition}}{}
\@ifundefined{lemma}{\newtheorem{lemma}{Lemma}}{}
\makeatother
\newcommand{\FF}{\mathbb{F}}
\newcommand{\beqn}{\begin{equation*}}
\newcommand{\beq}{\begin{equation}}
\newcommand{\eeqn}{\end{equation*}}
\newcommand{\eeq}{\end{equation}}
\title[Discreteness of asymptotic tensor ranks]{Discreteness of asymptotic tensor ranks}
\author{Jop Bri\"et, Matthias Christandl, Itai Leigh, Amir Shpilka, Jeroen Zuiddam}
\date{Source: arXiv:2306.01718v4 (CC BY 4.0)}
\begin{document}
\maketitle

\noindent\emph{Source and licence.} Discreteness of asymptotic tensor ranks. Version arXiv:2306.01718v4 (CC BY 4.0), distributed under the Creative Commons Attribution 4.0 International licence, as recorded in the arXiv licence metadata for this exact version and shown on the abstract page https://arxiv.org/abs/2306.01718v4. Full rights notice: licenses/arxiv-2306.01718-CC-BY-4.0.txt.\par\smallskip

\noindent\textit{Editorial scope.} Counted unit: the Lemma whose source label is \texttt{lem:upper-triangular} in the article (source0.tex lines 4426--4436, source label \texttt{lem:upper-triangular}), delivered together with the article's own complete original proof of it (source0.tex lines 4438--4510, one \texttt{proof} environment of 73 lines). No mathematics is written, repaired or re-derived: statement and proof are the article's own text line for line. Required context retained uncounted: def:pivot-matched (lines 4354--4367). Every definition, display and label of those lines is retained unchanged. Omitted from this package and not redistributed: 1--4353, 4368--4425, 4437--4437, 4511--4953. Nothing inside a retained excerpt is omitted or altered: every retained body is delivered exactly as the article's lines are written, so no omission receipt of any kind accompanies this package. Citations: the retained excerpts carry no citation command at all, so no bibliography entry of the article is transcribed and no reference is redistributed. Reference adaptation: none was needed and none was made; every reference inside the delivered unit resolves inside it, so no unresolved reference or citation is produced. Printed slips kept literal: no printed word, symbol, hypothesis or number of the counted statement or of its proof was altered. Layout and class: the article's own class is \texttt{daj} with packages including inputenc, mathtools, amsmath, amssymb, microtype, enumerate, accents, arydshln, amsthm, thmtools, thm-restate; the wrapper is the lane's pinned \texttt{amsart} editorial wrapper at 10pt on A4, which restates the theorem-like environments the retained text needs and adds the article's own macro definitions that the retained text needs (\texttt{\textbackslash FF}, \texttt{\textbackslash beq}, \texttt{\textbackslash beqn}, \texttt{\textbackslash eeq}, \texttt{\textbackslash eeqn}), with extra wrapper packages (bm, bbm, dsfont, xspace, cancel, mathtools). The delivered numbering is the unit's own. Assets: no retained span contains a figure environment, a graphics-inclusion command, a PDF-page inclusion, a table or a TikZ picture; no image file is copied into this package, the package declares no asset dependency and no image file is redistributed with it. Licence: version arXiv:2306.01718v4 of this article is distributed under the Creative Commons Attribution 4.0 International licence, as recorded in the arXiv licence metadata for this exact version and shown on the abstract page https://arxiv.org/abs/2306.01718v4. A full rights notice is delivered at licenses/arxiv-2306.01718-CC-BY-4.0.txt. Characters: every retained excerpt is pure ASCII, so the delivered engine, XeLaTeX with its default Latin Modern text font, reports no missing character.
\begin{definition}[Pivot-matched] \label{def:pivot-matched}
    We call $T\in\FF^{n}\otimes\FF^{n}\otimes\FF^{n}$ \emph{pivot-matched} if there are tensors $T_A$ and $T_B$ isomorphic 
    %\refnote{26}{What does "isomorphic" mean here?}\Jshort{Added it to the preliminaries} 
    to $T$ such that, 
    %
%
    letting $A_i = \sum_{j,k} (T_A)_{i,j,k}\, e_j \otimes e_k \in \FF^{n} \otimes \FF^{n}$ for $i \in [n]$ be the 1-slices of $T_A$, and $p_A=(f_A,g_A)$ the pivot map of these matrices, and
    letting $B_k = \sum_{i,j} (T_B)_{i,j,k}\, e_i \otimes e_j \in \FF^{n}\otimes \FF^{n}$ for $k \in [n]$ be the 3-slices of $T_B$, and $p_B=(f_B,g_B)$ the pivot map of these matrices, we have that $p_A$ and $p_B$ are injective and  for all $i \in [n]$ we have $f_A(i)=f_B(i)$.
%
%
%
%
%
\end{definition}

\begin{lemma}\label{lem:upper-triangular}
    Suppose $T\in\FF^n\otimes\FF^n\otimes\FF^n$ is a concise tensor that is %
    pivot-matched. 
    Then
    $T\boxtimes T$ restricts to a tensor of the form
\begin{multline}\label{eq:req-form}
\sum_{i=1}^n \big[e_i\otimes e_{f(i)}\big]\otimes\big[e_{f(i)}\otimes e_{g(i)}\big] \otimes \big[e_{h(i)}\otimes e_i\big]
+ \sum_{i=1}^n\sum_{j<k}\sum_{l,u,v=1}^n c_{ijkluv}\big[e_i\otimes e_j\big]\otimes \big[e_k\otimes e_l\big]\otimes \big[e_u\otimes e_v].
\end{multline}
for $f,g,h:[n]\to[n]$ functions such that $i\mapsto(f(i),g(i))$ is injective.
\end{lemma}

\begin{proof}
Let $A_1,\dots,A_n$ be the $1$-slices of $T_A$ as in \autoref{def:pivot-matched}, and $B_1,\dots,B_n$ the $3$-slices of $T_B$. Without loss of generality, we may assume that
\beq\label{eq:pivot1}
(A_i)_{p_A(j)} = (B_i)_{p_B(j)} = \delta_{ij}
\eeq
where $p_A=(f_A,g_A)$ is the pivot map for the $1$-slices of $T_A$ (rows enumerated by direction~$2$) and $p_B=(f_B,g_B)$ is for the $3$-slices of $T_B$ (rows by direction $1$), otherwise we can take invertible linear combinations of the slices to guarantee it.

%
%
%
%
Using these two slicings, we have the decomposition
\beqn
T^{\boxtimes 2}\geq T_A\boxtimes T_B
=
\sum_{i,u=1}^n\sum_{j,v=1}^n\sum_{k,w=1}^n (A_i)_{jk}(B_w)_{uv} (e_i\otimes e_u)\otimes (e_j\otimes e_v)\otimes (e_k\otimes e_w).
\eeqn

We apply the following three linear maps to the respective legs of this tensor:
\begin{align*}
U &= \sum_{i=1}^n E_{i,i}\otimes E_{f_A(i),f_A(i)}\\
V &= \sum_{v=1}^n E_{f_B(v),f_B(v)}\otimes E_{g_B(v),g_B(v)}\\
W &= \sum_{w=1}^nE_{g_A(w),g_A(w)}\otimes E_{w,w}.
\end{align*}
where $E_{i,i}$ is the projector onto $e_i$.
A small calculation then gives
\begin{multline*}
T^{\boxtimes 2}\geq(U\otimes V\otimes W)(T_A\boxtimes T_B)
=\\
\sum_{i,v,w=1}^n
(A_i)_{f_B(v)\,g_A(w)}\, (B_w)_{f_A(i)\,g_B(v)}\, 
(e_i\otimes e_{f_A(i)})\otimes (e_{f_B(v)}\otimes e_{g_B(v)})\otimes (e_{g_A(w)}\otimes e_w).
\end{multline*}
This is the final tensor. We now proceed to show how it fits the desired form \eqref{eq:req-form}.

By definition of $f_A$ as giving the first non-zero rows of the $A$ slices, we get non-zero contributions only from the $(i,v)$-pairs satisfying $f_B(v) \geq f_A(i)$ from the first factor after the triple sum.
This already shows that there is an upper-triangular structure present, by looking at the second component of the first leg and the first component of the second leg.

To finish showing $(U\otimes V\otimes W)(T_A\boxtimes T_B)$ is in the desired form we focus on the pairs $(i,v)$ satisfying $f_A(i) = f_B(v)$, aiming to show that this fits the first sum in \eqref{eq:req-form}. 
%\refnote{27}{Suggesting to number the equation at the statement of the lemma and referring to it here}\Jshort{Done}. %
Let~$S$ denote the set of such pairs.
Restricted to $S$, the triple sum becomes
\beqn
\sum_{(i,v)\in S}\sum_{w=1}^n 
(A_i)_{f_A(i)\,g_A(w)}\, (B_w)_{f_B(v)\,g_B(v)}\, 
(e_i\otimes e_{f_A(i)})\otimes (e_{f_B(v)}\otimes e_{g_B(v)})\otimes (e_{g_A(w)}\otimes e_w).
\eeqn
By~\eqref{eq:pivot1}, we have that 
\beqn
(B_w)_{f_B(v)\,g_B(v)} = (B_w)_{p_B(v)} = \delta_{vw}.
\eeqn
Hence, the sum over~$w$ collapses to a single term and we are left with the tensor 
\beqn
\sum_{(i,v)\in S}
(A_i)_{f_A(i)\,g_A(v)}\, 
(e_i\otimes e_{f_A(i)})\otimes (e_{f_B(v)}\otimes e_{g_B(v)})\otimes (e_{g_A(v)}\otimes e_v).
\eeqn
Fix $i,j$ such that $f_A(i) = j$.
Then~$v$ ranges over the set $f_B^{-1}(j)$, so all the $B$-slices whose pivot lies on row~$j$.
But by assumption, these are precisely the $A$-slices whose pivot lies on row~$j$.
This implies that the coordinate $(j,g_A(v))$ is the pivot of some $A$-slice.
Again by~\eqref{eq:pivot1} we  get that
\beqn
(A_i)_{f_A(i)\,g_A(v)} = \delta_{iv}.
\eeqn
Thus our tensor looks like
\beqn
\sum_{i=1}^n 
(e_i\otimes e_{f_A(i)})\otimes (e_{f_B(i)}\otimes e_{g_B(i)})\otimes (e_{g_A(i)}\otimes e_i).
\eeqn
Note that the maps $i \mapsto (i, f_A(i))$, $i \mapsto (f_B(i), g_B(i))$ and $i \to (g_A(i), i)$ are injective, as required
(i.e.\ this first sum is a diagonal of length~$n$).
\end{proof}

\end{document}
